\documentclass[10pt,a4paper]{amsart}
\usepackage[margin=19mm]{geometry}
\usepackage{amsmath,amssymb,mathtools}
\usepackage[scr=rsfs,cal=euler]{mathalfa}
\usepackage{xfrac}
\usepackage[unicode,hidelinks,bookmarks=false]{hyperref}
\newcommand{\ie}{i.e.\ }
\newcommand{\cf}{cf.\ }
\newcommand{\resp}{respectively\ }
\newcommand{\Subsection}{\subsection}
\providecommand{\nobreakdash}{\nobreak\hbox{-}}
\setlength{\emergencystretch}{2em}
\multlinegap0pt
\usepackage{float}
\usepackage{amssymb}
\usepackage{xcolor}
\usepackage[shortlabels]{enumitem}
\usepackage[normalem]{ulem}
\newtheorem{lemma}{Lemma}
\title{Nowhere-zero 8-flows in 3-edge-connected signed graphs}
\author{Matt DeVos, Kathryn Nurse, Robert \v{S}\'amal}
\date{Source: arXiv:2512.18923v1 (CC BY 4.0)}
\begin{document}
\maketitle

\noindent\emph{Source and licence.} Nowhere-zero 8-flows in 3-edge-connected signed graphs. Version arXiv:2512.18923v1 (CC BY 4.0), distributed by arXiv under the Creative Commons Attribution 4.0 International licence, http://creativecommons.org/licenses/by/4.0/, as recorded in the arXiv licence metadata for this exact version and shown on the abstract page https://arxiv.org/abs/2512.18923v1. Full licence notice: \texttt{licenses/arxiv-2512.18923-CC-BY-4.0.txt}.\par\smallskip

\noindent\textit{Editorial scope.} Counted unit: the article's lemma atMost2negCyc (source0.tex lines 385--387). The article's complete original proof of it is delivered with it (source0.tex lines 390--424, one \texttt{proof} environment of 35 lines). No mathematics is written, repaired or re-derived: the statement and the proof are the article's own lines, in the article's own notation, and no retained line is substituted or re-flowed.  Retained uncounted as the source-local context the proof needs: lines 370--372.  One declared adaptation: exactly 3 ancillary strings inside the retained spans are omitted (image-inclusion commands and, where the source repeats a label, the repeated label definitions) and no other line differs from the source; the figure environments, with their placement, captions and labels, are kept, and the package records the omitted strings in fidelity receipts bound to the affected bodies.  No retained line contains a citation, so the wrapper carries no bibliography and no bibliography entry.  The retained lines contain the article's own diagram code, which the wrapper typesets with the standard tikz packages; no image file is delivered and the package declares no asset dependency.  Editorial formatting only, in addition: an AMS \texttt{amsart} preamble that restates the article's own declarations and macros used by the retained lines, namely the environments \texttt{lemma} on separate counters, together with the standard packages those lines need.  The retained environments print in this document as Lemma 1 in order of appearance, whereas the article prints the same items with the numbering of its own full document.  The complete native source of the article as distributed in the arXiv e-print for this exact version is bundled unchanged as \texttt{sources/191433/source0.tex}.
\begin{lemma}\label{fragileNegCyc}
    Let $G$ be a 2-connected, fragile, well-behaved signed graph with no good theta-pair, and let $C \subseteq G$ be a good cycle. Then every negative cycle $N$ in $G - V(C)$ has $|V_3(G) \cap V(N)| = 2$ and $|V_2(G) \cap V(N)| \ge 1$. 
\end{lemma}

\begin{lemma}\label{atMost2negCyc}
    If $G$ is a $2$-connected, fragile, well-behaved signed graph with no good theta-pair, then for every good cycle $C$ the graph $G - V(C)$ has at most two negative cycles.
\end{lemma}

\begin{proof}
     Let $G' = G - E(C)$. By Lemma \ref{fragileNegCyc}, every negative cycle $D$ in $G'$ has $|V_3(G) \cap V(D)| = 2$ and $|V_2(G) \cap V(D)| \ge 1$.

    \begin{figure}[H]
        \centering
        
        %\caption{A component of $G-V(C)$ has two negative cycles.}
        \label{fig:f1}
    \end{figure}
    
    

    If a component $H$ of $G'$ contained two negative cycles $D,D'$, then there would be a good theta-pair using $D$ and a path $Q$ containing a degree-2 vertex from $D'$ and a degree 2 vertex from $C$, a contradiction. Hence every component of $G'$ has at most one negative cycle. 


    Now, suppose for contradiction that $G'$ has three negative cycles, say $D_1, D_2, D_3$. By 2-connectivity of $G$, for $1 \le i \le 3$ choose a pair of vertex disjoint paths from $V(D_i)$ to $V(C)$ and let $H_i$ be the subgraph of $G$ made by the union of $D_i$ and the two paths associated with it. Observe that by the previous paragraph, $H_1, H_2, H_3$ are pairwise vertex disjoint. For $1 \le i \le 3$ let $\{x_i,y_i\} = V(C) \cap V(H_i)$. Note that $H_i$ contains paths between $x_i$ and $y_i$ of both signs, and that at least one of those two paths contains a vertex in $V_2(G)$.

    \begin{figure}[H]
        \centering
        
        %\caption{An example where $x_i,y_i$ do not interleave $a,b$ on the left. An example where $x_i,y_i$ interleave $x_j,y_j$ on the right.}
        \label{fig:f34}
    \end{figure}

    Now, because $C$ is good, there exist $a,b \in V_2(G) \cap C$. It must be that, for each $1 \le i \le 3$, the vertices $x_i, y_i$ interleave $a,b$ on $C$, because otherwise there would be a good theta-pair using $D_i$ and a path $Q \subseteq H_i \cup C$ which contains $a,b$ as in the left side of the above figure. It also must be that $x_i,y_i$ do not interleave $x_j,y_j$ for $i \neq j$, because otherwise there would be a good theta-pair using $D_i$ and a path $Q\subseteq H_i \cup H_j \cup C$ containing a vertex in $V_2(G) \cap V(D_j)$, and a vertex in $V_2(G) \cap V(C)$, as in the right side of the above figure.

    \begin{figure}[H]
        \centering
        
        %\caption{An example where the vertices $a,x_1,x_2,x_3,b,y_3,y_2,y_1$ occur in that cyclic order around $C$}
        \label{fig:fragile5}
    \end{figure}

    Hence it follows (wlog) that the vertices $a,x_1,x_2,x_3,b,y_3,y_2,y_1$ occur in that cyclic order around $C$. But now $H_1 \cup C$ has an unbalanced theta while $H_2 \cup H_3 \cup C$ has a theta with a degree-2 vertex on each branch ($x_3$ and $y_3$ are the vertices of degree 3 in that theta), and so that these subgraphs are disjoint. Since the later subgraph contains a good cycle, this contradicts that $G$ is fragile and completes the proof.
\end{proof}

\end{document}
